\documentclass[conference]{IEEEtran}
\IEEEoverridecommandlockouts

\usepackage{cite}
\usepackage{amsmath,amssymb,amsfonts}
\usepackage{amsthm}
\usepackage{algorithmic}
\usepackage{graphicx}
\usepackage{textcomp}
\usepackage{xcolor}
\usepackage{booktabs}
\usepackage{array}
\usepackage{multirow}
\usepackage{url}
\usepackage{hyperref}
\usepackage{microtype}

\def\BibTeX{{\rm B\kern-.05em{\sc i\kern-.025em b}\kern-.08em
    T\kern-.1667em\kern-.125emE\kern-.125emX}}

% --- Theorem-like environments -----------------------------------------------
\newtheorem{theorem}{Theorem}
\newtheorem{proposition}[theorem]{Proposition}
\newtheorem{corollary}[theorem]{Corollary}
\newtheorem{lemma}[theorem]{Lemma}

\theoremstyle{definition}
\newtheorem{definition}{Definition}
\newtheorem{conjecture}[theorem]{Conjecture}

% --- Macros ------------------------------------------------------------------
\newcommand{\eg}{\emph{e.g.}}
\newcommand{\ie}{\emph{i.e.}}
\newcommand{\etal}{\emph{et~al.}}
\newcommand{\etc}{\emph{etc.}}
\newcommand{\src}[1]{\texttt{[#1]}}
\newcommand{\Bstar}{\ensuremath{B^{*}}}

\begin{document}

\title{A Navigation-Primitive Taxonomy of Forensic Evidence Sources:\\
Five Address Spaces and the Architecture of\\
Cross-Source Digital Investigation}

\author{\IEEEauthorblockN{Albert K.T. Hui}
\IEEEauthorblockA{\textit{Security Ronin} \\
albert@securityronin.com}}

\maketitle

\begin{abstract}
Digital forensic frameworks classify evidence by \emph{artifact type} -- what
the evidence \emph{is} -- not by the mechanism required to address and
traverse it. This distinction is consequential: artifact-type classification
does not determine layer dependencies, parser reusability, or
epistemological properties such as attacker-durability. We propose to
classify evidence sources by \emph{navigation primitive} -- the abstract
addressing scheme by which the bytes of evidence are reached -- and show
that exactly five primitives are sufficient to cover contemporary forensic
practice: persistent storage \src{P}, memory \src{M}, log \src{L}, live
query \src{Q}, and content-addressed \src{C}. We formalize evidence sources
as triples $(A, f_{\text{nav}}, f_{\text{parse}})$ where $f_{\text{nav}}: A
\rightharpoonup B^{*}$ is a navigation function valued in byte sequences and
$f_{\text{parse}}: B^{*} \to R$ is a parser valued in records. From this
definition we prove a \emph{Parser Independence} theorem: any two evidence
sources with the same parser produce identical records on identical byte
inputs, regardless of the navigation primitive that retrieved them; we
further establish that the dependency rule (PARSER must not import NAVIGATE)
is the contrapositive of this theorem and is thus a formal architectural
constraint. We define \emph{attacker-durability} of captured evidence and
prove that \src{Q} achieves durability by transmission to investigator
storage (Proposition~\ref{prop:q-durability}) and that \src{C} achieves
durability \emph{intrinsically} from collision-resistance of the addressing
hash (Proposition~\ref{prop:c-durability}); we prove as a corollary that
\src{C} is the unique source type whose durability requires no external
trust anchor. We give a schema-conscious definition of the
\emph{hash-pivot} operation that yields a typed-multiset join across CAS
stores. The taxonomy is hash-agnostic: SHA-1, SHA-256, and future
hash-function migrations are accommodated by a single collision-resistance
parameter. We describe \emph{Issen}, a Rust reference implementation
organized as thirteen crates across five layers (KNOWLEDGE, CONTAINER,
NAVIGATE, PARSER, ORCHESTRATION) whose dependency rules are mechanically
forced by the taxonomy. Four case studies evaluate the framework: a
rootkit-concealed cryptominer requiring \src{P}+\src{M}+\src{Q} fusion; a
supply-chain compromise resolved by hash-pivot through \src{C} stores; an
empirical Parser Independence demonstration with byte-identical
SHA-256 hashes of parser output across four navigation paths into the same
Chrome \texttt{History} file; and a counterfactual case showing concretely
how a \src{P}+\src{M}+\src{L}-only triage tool misses an
\texttt{ld.so.preload} race that \src{Q} captures. The taxonomy supplies
the missing formal foundation for cross-source correlation called for in
Garfinkel's 2010 agenda.
\end{abstract}

\begin{IEEEkeywords}
digital forensics, evidence taxonomy, memory forensics, content-addressed
storage, supply chain forensics, live response, navigation primitive,
parser independence
\end{IEEEkeywords}

\section{Introduction}
\label{sec:introduction}

A forensic analyst sits before three tools and one question. The question is
simple: what URLs did the user visit on this machine? The disk image is open
in one tool, which presents a tabular view of Chrome \texttt{History} entries
extracted from the SQLite database living on the NTFS volume. The memory
dump is open in a second tool, which surfaces fragments of browser-related
strings recovered from carved pages. A third tool displays a Velociraptor
collection that pulled the live \texttt{History} file moments before the
endpoint was isolated. Three tools, three incompatible record schemas, three
different timestamp conventions, three different notions of what a ``URL
visit'' even is. The underlying artifact -- the Chromium \texttt{urls} table
and its \texttt{visits} sibling -- is identical across all three.

The concrete consequence is not abstract: the analyst must (a) write a
bespoke per-source adapter to reconcile the schemas (a multi-day engineering
task, performed under investigative time pressure, frequently re-done from
scratch in the next case), or (b) re-parse the artifact manually for each
source from raw bytes (which discards the parsing infrastructure of all
three tools), or (c) accept that evidence from one source is forensically
siloed from the others, weakening the case. None of these is acceptable; all
three occur in practice. The analyst's frustration is not a tooling
accident: it is the predictable consequence of an architectural choice that
the field has internalized without examining: \emph{evidence is classified
by what it is, not by how it is reached.}

Existing taxonomies of digital evidence are, almost without exception,
artifact-type taxonomies. Casey's textbook organizes evidence by category
(operating system, application, network, mobile)~\cite{casey2011digital}.
The CASE/UCO ontology models artifacts as typed
nodes~\cite{caseuco2018,casey2017advancing}. Carrier's data-abstraction
model layers physical, file system, and application
views~\cite{carrier2003getting,carrier2005file}. McKemmish's foundational
definition partitions forensic computing into identification, preservation,
analysis, and presentation phases against artifact
classes~\cite{mckemmish1999forensic}. Each is valuable, and each is
\emph{orthogonal} to the architectural question that determines whether the
analyst above can ever get a unified answer.

The architectural question is the address space. Disk evidence is reached by
walking a filesystem namespace from name to inode to block. Memory evidence
is reached by walking a process list, then a page table from virtual to
physical address. Log evidence is reached by seeking by timestamp or
sequence number through chunked records. Live-query evidence is reached by
dispatching a query to an endpoint and paginating result rows.
Content-addressed evidence is reached by following hashes through a Merkle
DAG. These are not five flavors of the same operation; they are five
fundamentally different ways of locating bytes -- each with distinct
dependency, trust, and durability properties. Table~\ref{tab:glance}
gives the map.

\begin{table}[h]
\centering\small
\begin{tabular}{@{}p{0.06\columnwidth}p{0.38\columnwidth}p{0.44\columnwidth}@{}}
\toprule
& \textbf{Navigation step} & \textbf{Key property} \\
\midrule
\src{P} & name $\to$ inode $\to$ block & durable by acquisition \\
\src{M} & PID $\to$ EPROCESS $\to$ VA $\to$ PA & volatile; instantaneous snapshot \\
\src{L} & timestamp $\to$ record boundary & append-only; durable by acquisition \\
\src{Q} & (endpoint, query) $\to$ result rows & produced, not stored; durable by transmission \\
\src{C} & hash $\to$ blob $\to$ Merkle DAG & immutable; intrinsically durable \\
\bottomrule
\end{tabular}
\caption{Five navigation primitives at a glance.}
\label{tab:glance}
\end{table}

We propose a taxonomy of five \emph{navigation primitives} and show that
classification by navigation primitive yields properties that artifact-type
classification cannot: parser reusability, attacker-durability profiles, and
a layered architecture whose dependency rules are mechanically enforceable.
We instantiate the framework in \emph{Issen}, a Rust platform of thirteen
crates whose layer membership is derived from the taxonomy.

\paragraph{Contributions.}
We make five claims, each tied to the section that delivers it.
\begin{enumerate}
    \item \textbf{Contribution 1 (Section~\ref{sec:taxonomy}).} The
    \emph{navigation-primitive} criterion for source classification, formalized
    as an equivalence relation $\sim$ over address-space algebras
    (Definition~\ref{def:nav-primitive}). The taxonomically consequential
    property of an evidence source is its addressing scheme, not its artifact
    category. This yields exactly five primitives under the criterion --
    \src{P}, \src{M}, \src{L}, \src{Q}, \src{C} -- defined by distinct
    address-space algebras. The three-source folk model (P/M/L) is the proper
    subset that emerges when one ignores addressing.

    \item \textbf{Contribution 2 (Section~\ref{sec:properties}, Proposition~\ref{thm:parser-indep}).}
    The \emph{Parser Independence} principle: identical bytes yield identical
    records under a fixed deterministic parser, regardless of the navigation
    primitive that retrieved the bytes. Naming it converts a folk design
    principle into an explicitly testable claim -- demonstrated empirically
    with byte-identical SHA-256 hashes across four navigation paths
    (Section~\ref{sec:case3}).

    \item \textbf{Contribution 3 (Section~\ref{sec:properties}, Propositions~\ref{prop:q-durability} and~\ref{prop:c-durability}, Corollary~\ref{cor:c-anchor-free}).}
    Formalization of \src{Q} as a first-class source type with attacker-durability
    by transmission and of \src{C} as a first-class source type with intrinsic
    attacker-durability; we prove that \src{C} is the unique source type whose
    durability requires no external trust anchor. To our knowledge, neither
    \src{Q} nor \src{C} has previously been admitted to the forensic source
    taxonomy as a peer of the disk-memory-log triad.

    \item \textbf{Contribution 4 (Section~\ref{sec:properties}, Proposition~\ref{thm:layer-rule}, and Section~\ref{sec:implementation}).}
    The \emph{Layer Dependency Rule}: a parser that imports a navigation
    component cannot be source-agnostic, because its domain is then
    $A_\pi$, not $\Bstar$. This is the contrapositive of Parser Independence.
    Issen instantiates the rule in Cargo crate boundaries enforced by the
    build system.

    \item \textbf{Contribution 5 (Sections~\ref{sec:cases}~and~\ref{sec:discussion}).}
    Four case studies, including (i) an empirical Parser Independence
    demonstration with byte-identical SHA-256 hashes across four sources,
    (ii) a counterfactual case showing concretely how a \src{P}+\src{M}+\src{L}-only
    triage tool fails to capture evidence that \src{Q} captures, and (iii)
    an explicit completeness argument framed as an open conjecture rather
    than mere absence of counterexamples.
\end{enumerate}


\section{Background}
\label{sec:bg}

\subsection{Process Models and Their Silences}

Digital forensic process models describe how an investigation should
proceed, not what classes of evidence exist. The DFRWS framework first
codified the seven-phase pipeline of identification, preservation,
collection, examination, analysis, presentation, and decision~\cite{palmer2001roadmap}.
Carrier and Spafford generalized this into a five-phase \emph{Integrated
Digital Investigation Process} that explicitly aligns the digital phases
with the physical-crime-scene phases that
preceded them~\cite{carrier2003getting}. Beebe and Clark extended the
emphasis on hierarchical objectives, arguing that procedural decomposition
must respond to investigative
objectives rather than artifact taxonomy~\cite{beebe2005hierarchical}.
Carrier's later dissertation~\cite{carrier2006hypothesis} formalized
hypothesis testing as the epistemic core of digital investigation: every
analytical claim is a hypothesis that must be falsifiable from the data.

These models are silent on the question we ask: \emph{what kinds of addressing
does a piece of evidence support, and what are the architectural consequences?}

\subsection{The Three-Source Folk Model and Its Failure Modes}

The implicit, never-formalized folk model of forensic evidence partitions
sources into three: \emph{disk}, \emph{memory}, and
\emph{log}~\cite{kent2006sp80086,casey2011digital,carrier2005file}. Standard
practitioner training and toolchains -- The Sleuth Kit and
Autopsy~\cite{carrier2005file}, Volatility for memory analysis~\cite{ligh2014artmemory,case2017memory},
and timeline tools such as Plaso/log2timeline~\cite{plaso} -- align with this
partition. The folk model has served the field well for two decades but is
breaking under the weight of two developments.

The first is the rise of \emph{enterprise live response} and endpoint
detection-and-response telemetry. When a Velociraptor agent~\cite{velociraptor}
runs a VQL query and returns a result set of currently-open TCP connections,
that result is not on a disk image; it is not yet in any log; it lives,
briefly, only in the network packet that delivers it from the endpoint to
the investigator. Cohen, Bilby, and Garfinkel argued in
2011~\cite{cohen2011distributed} that this kind of capture is a fundamentally
different forensic operation, but the field continued to treat the resulting
data as ``just another log'' -- an analytical convenience that obscures the
property that the data did not exist before the query was issued and, once
captured and transmitted, cannot be modified by an adversary regardless of
subsequent on-host actions.

The second development is \emph{content-addressed} ecosystems. Git,
Sigstore's transparency log~\cite{newman2022sigstore}, in-toto~\cite{torresarias2019intoto},
OCI image manifests~\cite{ociimagespec}, and IPFS~\cite{benet2014ipfs} all
address artifacts by the cryptographic hash of their content; references between
artifacts are themselves hashes. The resulting Merkle DAG has no path, no virtual
address, no timestamp ordering -- only hash $\to$ blob $\to$ content graph.
Supply-chain attacks (SolarWinds~\cite{cisa2021solarwinds},
\texttt{xz-utils}~\cite{cve2024xz}) forced the forensic community to reckon
with evidence trails that live entirely in content-addressed stores that no
traditional triage tool addresses.

\subsection{The Gap}

The literature provides three separate strands -- traditional disk and log
forensics, memory forensics as it matured through Volatility and
Rekall~\cite{ligh2014artmemory,case2017memory}, and emerging work in live
response and supply-chain forensics -- that have not been unified by a
common formal model.

The closest representation-level effort is the CASE/UCO
ontology~\cite{caseuco2018,casey2017advancing}, which provides a typed
vocabulary for describing digital evidence \emph{after} it has been parsed.
Our taxonomy is complementary, not competing: CASE/UCO is \emph{annotation}
(the schema in which findings are exchanged), while the navigation-primitive
taxonomy is \emph{architecture} (the addressing discipline by which findings
are produced). A CASE/UCO record consumed by a downstream tool conveys the
\emph{what}; a navigation-primitive tag (\eg{}, \texttt{nav-primitive: Q})
conveys the \emph{how}, which CASE/UCO does not encode and which determines
the trust profile and durability claims that the consumer can make.

The closest infrastructure-level effort is Carrier's data-abstraction model
and its concrete instantiation in The Sleuth Kit (TSK)~\cite{carrier2005file,tsk_repo}.
TSK exposes \texttt{img\_*}, \texttt{vol\_*}, \texttt{fs\_*}, and \texttt{file\_*}
APIs corresponding to a layered traversal: image $\to$ volume $\to$
filesystem $\to$ file. This layering is foundational for \src{P} navigation
specifically: TSK's \texttt{img\_} layer is what we generalize as CONTAINER,
and TSK's \texttt{fs\_}/\texttt{file\_} layers together constitute the
\src{P}-specialization of NAVIGATE in our framework. The crucial difference
is that TSK's layering is \emph{specific} to disk evidence -- it has no
analog of \texttt{img\_} for memory dumps, no analog of \texttt{vol\_} for
log streams, no analog of \texttt{file\_} for hash-addressed blobs. Our
taxonomy generalizes Carrier's layered insight to four other addressing
schemes that did not exist in their current form when TSK was designed.

Garfinkel's 2009 work on automating disk forensic processing with TSK, XML,
and Python~\cite{garfinkel2009automating} provides early evidence that even
within \src{P} alone the absence of a uniform addressing model forces
ad-hoc adapters: Garfinkel's DFXML output exists precisely because tool-output
schemas were not interoperable. The same problem at five-source scale
motivates the present paper.

Garfinkel's 2010 agenda~\cite{garfinkel2010digital} identified cross-source
correlation as one of the key challenges of the coming decade; sixteen
years later, the field has rich tooling for each source individually but no
formal account of what ``cross-source'' means. We aim to provide that
account.

\section{The Navigation-Primitive Taxonomy}
\label{sec:taxonomy}

Table~\ref{tab:comparison} shows the five primitives. In brief:
\src{P} walks a filesystem tree (name~$\to$~inode~$\to$~block);
\src{M} walks a kernel object graph (PID~$\to$~EPROCESS~$\to$~VA~$\to$~PA);
\src{L} seeks a record stream by timestamp or cursor;
\src{Q} dispatches a query to an endpoint and receives result rows that did
not exist before the query was issued;
\src{C} follows cryptographic hashes through a Merkle DAG where the address
\emph{is} the hash of the content.
The formal definitions below make these distinctions precise.

\begin{table*}[!t]
\centering
\caption{The Five Navigation Primitives}
\label{tab:comparison}
\begin{tabular}{@{}p{0.06\textwidth}p{0.18\textwidth}p{0.20\textwidth}p{0.20\textwidth}p{0.10\textwidth}p{0.14\textwidth}@{}}
\toprule
\textbf{Type} & \textbf{Address space} & \textbf{Traversal operation} & \textbf{Canonical tools} & \textbf{Mutability} & \textbf{Attacker durability} \\
\midrule
\src{P} & path / inode / block & directory walk; inode lookup; block read & TSK, Autopsy, ewf, ext4fs-forensic & mutable on host & by acquisition \\
\src{M} & (PID, VA) $\to$ PA & process-list walk; page-table walk & Volatility, Rekall, memf-* & mutable on host (DKOM) & by acquisition \\
\src{L} & timestamp / cursor & seek to record boundary; field decode & Plaso, EvtxECmd, winevt-forensic & append-only by intent & by acquisition \\
\src{Q} & (endpoint, query, cursor) & query dispatch; cursor pagination & Velociraptor, OSQuery, WMI & N/A (produced, not stored) & by transmission \\
\src{C} & cryptographic hash & hash $\to$ blob; reference following & git, Sigstore, OCI, IPFS & immutable by construction & intrinsic \\
\bottomrule
\end{tabular}
\end{table*}

\subsection{Formal Foundations}

We treat byte sequences, denoted $\Bstar = \{0,1,\ldots,255\}^{*}$, as the
unstructured medium through which all forensic evidence ultimately passes.
Forensic records, denoted $R$, are the typed objects that populate analyst
reports, timelines, and ontologies (\eg{}, browser history rows, EVTX event
records, EPROCESS structures). Parsers translate from $\Bstar$ to $R$.
Navigation translates from addresses to $\Bstar$. We work in the category
$\mathbf{PFn}$ of sets and partial functions~\cite{maclane1998categories}:
objects are sets, morphisms are partial functions, composition is the usual
chain composition with the convention that an undefined intermediate value
propagates as undefined.

\begin{definition}[Evidence Source]
\label{def:source}
An \emph{evidence source} is a triple $S = (A,\, f_{\mathrm{nav}},\,
f_{\mathrm{parse}})$, where:
\begin{itemize}
    \item $A$ is an \emph{address space}: a set whose elements name reachable
    locations within the source;
    \item $f_{\mathrm{nav}} : A \rightharpoonup \Bstar$ is a (partial)
    \emph{navigation function} that, given an address, yields the byte
    sequence stored at that address; and
    \item $f_{\mathrm{parse}} : \Bstar \to R$ is a \emph{parser} that
    interprets bytes as forensic records.
\end{itemize}
The composition $f_{\mathrm{parse}} \circ f_{\mathrm{nav}} : A \rightharpoonup R$
yields records keyed by source address.
\end{definition}

\begin{definition}[Address-Space Algebra]
\label{def:addr-algebra}
An \emph{address-space algebra} is a tuple $\mathcal{A} = (A, \mathrm{Op}, \to)$
where $A$ is a set of addresses, $\mathrm{Op}$ is a finite signature of
\emph{traversal operations} (each $\omega \in \mathrm{Op}$ is a partial
function $\omega : A^{k_\omega} \to A$ for some arity $k_\omega \geq 0$),
and $\to \;\subseteq A \times A$ is the reachability relation generated by
$\mathrm{Op}$ (\ie{}, $a \to a'$ iff some $\omega \in \mathrm{Op}$ maps $a$
to $a'$, and $\to$ is taken under transitive-reflexive closure).
\end{definition}

\begin{definition}[Navigation Primitive]
\label{def:nav-primitive}
A \emph{navigation primitive} $\pi$ is an isomorphism class of address-space
algebras under the equivalence relation $\sim$ defined as follows. Two
algebras $\mathcal{A} = (A, \mathrm{Op}, \to)$ and $\mathcal{A}' = (A',
\mathrm{Op}', \to')$ are equivalent, $\mathcal{A} \sim \mathcal{A}'$, iff
there exist a bijection $\varphi : A \to A'$ and a one-to-one
correspondence $\sigma : \mathrm{Op} \to \mathrm{Op}'$ that preserves arity
and that satisfies $\varphi(\omega(a_1, \ldots, a_{k_\omega})) =
\sigma(\omega)(\varphi(a_1), \ldots, \varphi(a_{k_\omega}))$ wherever the
left-hand side is defined. The bijection condition guarantees no
information loss; the operation correspondence guarantees that the same
traversal questions can be answered in either algebra.
\end{definition}

\begin{definition}[Source Type]
\label{def:source-type}
Two evidence sources $S = (A, f_{\mathrm{nav}}, f_{\mathrm{parse}})$ and
$S' = (A', f'_{\mathrm{nav}}, f'_{\mathrm{parse}})$ are of the same
\emph{source type} iff their underlying address-space algebras are equivalent
under $\sim$ (Definition~\ref{def:nav-primitive}), \ie{}, they share a
navigation primitive $\pi$.
\end{definition}

Two sources are the same type iff their addressing question -- ``starting at
$a$, where can I go and by what operation?'' -- is substitutable. An EVTX
file and a syslog stream are both \src{L}. A Chrome \texttt{History} file on
a mounted disk (\src{P}) and the same file recovered from carved memory (\src{M})
are not the same type, despite being the same artifact.

\subsection{The Five Primitives}

\subsubsection{\src{P} -- Persistent Storage}

\paragraph{Address space.} The filesystem namespace, $A_{P} = \mathrm{Path}
\cup \mathrm{Inode} \cup \mathrm{Block}$, with structure-preserving maps
$\mathrm{Path} \to \mathrm{Inode} \to \mathrm{Block}^{*}$ supplied by the
filesystem.
\paragraph{Traversal.} Resolve a path to an inode by directory walk; resolve
the inode to a block list by inode-table lookup; resolve the block list to
bytes by issuing reads against the underlying container (sector stream from
an E01 image, raw partition, FUSE mount, \etc{}).
\paragraph{Canonical examples.} E01/EWF disk images~\cite{ewf2011}, raw
partitions, NTFS, ext4, APFS, FUSE-bridged images.
\paragraph{Epistemological properties.} High completeness when the volume is
acquired offline; vulnerable to anti-forensic deletion and timestomping;
records are mutable in principle but typically frozen by acquisition.

\subsubsection{\src{M} -- Memory}

\paragraph{Address space.} Per-process virtual address space $A_{M} =
\mathrm{PID} \times \mathrm{VA}$, layered over the OS object graph and
ultimately resolving through page tables to physical addresses
$\mathrm{PA}$.
\paragraph{Traversal.} Walk the kernel's process list to reach an
\texttt{EPROCESS} (Windows) or \texttt{task\_struct} (Linux); walk the VAD
tree (Windows) or VMA list (Linux) to enumerate mapped regions; walk page
tables (PML4 on x86\_64, PAE on legacy x86, the four-level scheme on
AArch64) to map a virtual address to a physical
page~\cite{ligh2014artmemory,case2017memory}.
\paragraph{Canonical examples.} WinPMEM, LiME, AVML, hiberfil.sys,
Windows crash dumps.
\paragraph{Epistemological properties.} Captures non-persistent state
unavailable from \src{P}; vulnerable to direct kernel object manipulation
(DKOM); contents are an instantaneous snapshot rather than a trace.

\subsubsection{\src{L} -- Log}

\paragraph{Address space.} A linearly ordered space of cursors, $A_{L}
\subseteq \mathbb{N} \cup \mathrm{Timestamp} \cup \mathrm{Cursor}$, where a
cursor names a record boundary within a stream.
\paragraph{Traversal.} Seek to a chunk or block boundary by timestamp or
sequence number; read forward, decoding record boundaries; decode each record
to a structured field set. Backward traversal requires either an index or
linear scan.
\paragraph{Canonical examples.} Windows EVTX binary
format~\cite{microsoft2023evtx}, systemd-journald binary
journal~\cite{poettering2012journal}, Apple Unified Logging
\texttt{tracev3}~\cite{apple2016unifiedlogging}, AWS CloudTrail JSON event
streams~\cite{aws2024cloudtrail}, PCAP-NG capture
files~\cite{pcapng2024}.
\paragraph{Epistemological properties.} Append-only by intent (mutation
without trace requires privileged on-host action); linearly ordered, which
greatly simplifies temporal correlation; subject to rollover and rotation
that can lose records.

\subsubsection{\src{Q} -- Live Query}

\paragraph{Address space.} The Cartesian product of well-formed query
expressions and pagination cursors, $A_{Q} = \mathrm{Endpoint} \times
\mathrm{Query} \times \mathrm{Cursor}$.
\paragraph{Traversal.} Dispatch a query to an endpoint; the endpoint
executes the query against live state; results stream back as typed rows,
with cursor tokens advancing pagination. Re-execution of the same query at a
different time generally yields a different result.
\paragraph{Canonical examples.} Velociraptor VQL~\cite{velociraptor},
WMI/WQL~\cite{microsoft2023wmi}, OSQuery SQL~\cite{osquery}, GRR
Rapid Response~\cite{cohen2011distributed}.
\paragraph{Epistemological properties.} Evidence is \emph{produced} by the
act of querying; it does not pre-exist in any addressable store. Once
captured to an investigator-controlled system, the result is durable against
on-host adversarial action (Section~\ref{sec:properties}). The query itself
is part of the evidence chain and must be recorded.

\subsubsection{\src{C} -- Content-Addressed}

\paragraph{Address space.} The image of a cryptographic hash function, $A_{C}
\subseteq \{0,1\}^{n}$ for hash output width $n$ (\eg{}, SHA-1, SHA-256). The
\emph{address is the hash}.
\paragraph{Traversal.} Given a hash, retrieve the blob whose hash is that
value; parse the blob's references (each itself a hash); follow references
to traverse a Merkle DAG. Cycles are impossible because each node's hash
depends on the hashes of its references.
\paragraph{Canonical examples.} Git commit/tree/blob graphs~\cite{chacon2014progit},
OCI image manifests and layers~\cite{ociimagespec}, IPFS~\cite{benet2014ipfs},
the Sigstore Rekor transparency log~\cite{newman2022sigstore},
in-toto links~\cite{torresarias2019intoto}.
\paragraph{Epistemological properties.} Identity equals content. Modification
yields a new address rather than altering the artifact at the original
address (Section~\ref{sec:properties}). Enables \emph{hash-pivot}:
simultaneous query of all known CAS stores by a single hash.

\subsection{Layer Hierarchy}

The taxonomy yields a five-layer architecture in which each layer has a
defined dependency relationship with the layers below. Figure~\ref{fig:layers}
sketches the dependencies; Section~\ref{sec:implementation} maps these
layers to concrete crates in the Issen reference implementation.

\begin{figure}[t]
\centering
\begin{tabular}{@{}l@{}}
\toprule
\textbf{ORCHESTRATION}~~ wires sources, correlates evidence, drives UX \\
~~$\uparrow$ depends on all layers below \\
\midrule
\textbf{PARSER}~~~~~~~~~~~ $\Bstar \to R$, source-agnostic \\
~~$\uparrow$ depends only on KNOWLEDGE \\
\midrule
\textbf{NAVIGATE}~~~~~~~~~ $A \to \Bstar$, primitive-specific\\
~~ \src{P}: filesystem ~~~~~ \src{M}: paging / OS structure \\
~~ \src{L}: log format ~~~~~~ \src{Q}: query engine \\
~~ \src{C}: graph navigation \\
~~$\uparrow$ depends on CONTAINER + KNOWLEDGE \\
\midrule
\textbf{CONTAINER}~~~~~~ raw stream from outer format (E01, dump, \ldots) \\
~~$\uparrow$ depends on KNOWLEDGE \\
\midrule
\textbf{KNOWLEDGE}~~~~ format constants, magics, schemas \\
\bottomrule
\end{tabular}
\caption{Five-layer architecture induced by the taxonomy. Higher layers
depend only on lower layers. PARSER is independent of NAVIGATE because both
operate at the byte-sequence interface (Proposition~\ref{thm:parser-indep}).}
\label{fig:layers}
\end{figure}

The KNOWLEDGE layer holds purely declarative facts about formats: page
sizes, magic numbers, struct field offsets, schema definitions. The
CONTAINER layer decodes outer dump or image formats (E01 segments, WinPMEM
output, ELF core dumps) into a raw addressable stream; \src{Q} and \src{C}
have no separate container layer because the endpoint or hash store \emph{is}
the entry point. The NAVIGATE layer implements a navigation primitive
against a stream or store. The PARSER layer interprets bytes as records and
is, by construction, independent of which navigation primitive produced
those bytes. The ORCHESTRATION layer is the only place that may depend on
all of the above, and is responsible for routing artifacts to parsers and
for cross-source correlation.

We will see in Section~\ref{sec:properties} that the dependency rules are not
merely a stylistic convention; they are forced by Proposition~\ref{thm:parser-indep}
and by the contrapositive of Proposition~\ref{thm:layer-rule}.

\section{Formal Properties}
\label{sec:properties}

This section delivers Contributions~2,~3, and~4. We prove a semantically
strong Parser Independence theorem, derive the layer dependency rule from
its contrapositive, formalize attacker-durability, characterize the distinct
durability profiles of \src{Q} and \src{C}, and give a schema-conscious
definition of hash-pivot.

\subsection{Parser Independence}

The theorem we want is not merely a type-checking observation. We need a
\emph{semantic} claim: that two distinct navigation paths into the
``same'' artifact yield the \emph{same} record under a fixed parser. We
make this precise with the following theorem, which has two halves -- a
syntactic well-typedness half, and a semantic equality half.

\begin{theorem}[Parser Independence]
\label{thm:parser-indep}
Let $p : \Bstar \to R$ be a deterministic, side-effect-free parser, and let
$S = (A, f_{\mathrm{nav}}, p)$ and $S' = (A', f'_{\mathrm{nav}}, p)$ be two
evidence sources with possibly distinct navigation primitives $\pi$ and
$\pi'$. Then:

\smallskip
\noindent\textbf{(I) Type composition.}
Both compositions $p \circ f_{\mathrm{nav}}: A \rightharpoonup R$ and $p
\circ f'_{\mathrm{nav}}: A' \rightharpoonup R$ are well-typed and yield
records in $R$.

\smallskip
\noindent\textbf{(II) Semantic invariance.}
For any addresses $a \in A$ and $a' \in A'$ with $f_{\mathrm{nav}}(a) =
f'_{\mathrm{nav}}(a') = b \in \Bstar$ (\ie{}, the two navigation paths
return identical bytes), $p(f_{\mathrm{nav}}(a)) = p(f'_{\mathrm{nav}}(a'))$
as values in $R$. The equality holds bit-for-bit under any canonical
serialization of $R$.
\end{theorem}

\begin{proof}
\textbf{(I).} By Definition~\ref{def:source} the parser's domain is $\Bstar$,
independent of $A$. Both $f_{\mathrm{nav}}: A \rightharpoonup \Bstar$ and
$f'_{\mathrm{nav}}: A' \rightharpoonup \Bstar$ have codomain $\Bstar$,
matching the domain of $p$. Composition in $\mathbf{PFn}$ is defined when
the codomain of the inner function matches the domain of the outer; thus
$p \circ f_{\mathrm{nav}}$ and $p \circ f'_{\mathrm{nav}}$ are both
well-typed in $\mathbf{PFn}$. The construction does not reference $\pi$ or
$\pi'$.

\smallskip\noindent
\textbf{(II).} By assumption, $p$ is deterministic and side-effect-free: $p$
is mathematically a total function on its domain modulo the partiality
inherited from $f_{\mathrm{nav}}$, with no implicit dependency on global
state. Determinism implies $p(b) = p(b)$ for any fixed $b \in \Bstar$.
Substituting $b = f_{\mathrm{nav}}(a) = f'_{\mathrm{nav}}(a')$ yields
$p(f_{\mathrm{nav}}(a)) = p(b) = p(f'_{\mathrm{nav}}(a'))$. Bit-for-bit
equality of canonical serializations follows because record canonicalization
is itself a deterministic function of the record. \qedhere
\end{proof}

Half (I) is the syntactic content -- type composition -- which has been an
implicit design principle in tools like TSK and Volatility but has not
previously been stated formally for cross-source composition. Half (II) is
the semantic content and is the load-bearing claim of the theorem: it says
that for a properly written parser, four distinct navigation paths into
the ``same'' bytes (\src{P}-live, \src{P}-image, \src{M}-carve, \src{Q}-pull
of the Chrome \texttt{History} file in Section~\ref{sec:case3}) produce
identical records, and -- because canonical serializations are
deterministic -- identical SHA-256 hashes of those records. This converts
Parser Independence from a folk principle into an empirically falsifiable
prediction.

\paragraph{Scope of the determinism assumption.} The theorem requires the
parser to be free of nondeterminism (\eg{}, hash randomization of internal
maps, system-clock-dependent fields) and free of externalized state (\eg{},
no calls to schema-resolution services). These are standard
``pure-function'' constraints; we discuss empirical compliance as a
threat-to-validity (Section~\ref{sec:discussion}).

\subsection{Layer Dependency Rule}

The dependency rule is best understood not as an independent theorem but as
the contrapositive of Theorem~\ref{thm:parser-indep}. We state it
correspondingly. The structure is: \emph{if Parser Independence is required
of a parser, then the parser must not depend on NAVIGATE.}

\begin{theorem}[Layer Dependency Rule, Contrapositive Form]
\label{thm:layer-rule}
Let $p$ be a parser, and let $\Pi$ denote the set of navigation primitives
that some deployment of $p$ may encounter. If $|\Pi| \geq 2$ and $p$ is
required to satisfy Parser Independence (Theorem~\ref{thm:parser-indep})
across $\Pi$, then $p$'s domain must equal $\Bstar$ -- equivalently, $p$
must not import or otherwise depend on any NAVIGATE component specific to a
particular $\pi \in \Pi$.
\end{theorem}

\begin{proof}
Theorem~\ref{thm:parser-indep} states that semantic invariance holds for
parsers $p : \Bstar \to R$. The hypothesis is unconditional in $\pi$: the
domain of $p$ is $\Bstar$, not any $A_\pi$. Now suppose $p$ has domain
$A_\pi$ for some particular $\pi$ -- equivalently, $p$ depends on a
NAVIGATE component fixing the address space to $A_\pi$. Then $p$ is not a
function on $\Bstar$, and Theorem~\ref{thm:parser-indep} does not apply to
$p$. We must check that $p$ also \emph{cannot} satisfy Parser Independence
under any alternative argument. Take $\pi' \neq \pi$ in $\Pi$ and a source
$S' = (A', f'_{\mathrm{nav}}, p)$ of primitive $\pi'$. The domain of $p$
contains addresses of type $A_\pi$ but not $A_\pi$-typed values produced
through $f'_{\mathrm{nav}}$, which produces values in $\Bstar$ (or in
$A'$). The composition $p \circ f'_{\mathrm{nav}}$ is therefore not even
well-typed. Hence $p$ fails Type composition (Theorem~\ref{thm:parser-indep}.I)
and \emph{a fortiori} fails Semantic invariance. The contrapositive
direction is established without circularity: we have shown that the
\emph{absence} of an unconditional-in-$\pi$ domain forces a type failure on
some pair of primitives, not that a parser ``cannot be reused'' as a
premise. \qedhere
\end{proof}

The contrapositive is the practitioner's rule of thumb: \emph{if your
parser imports a navigation crate, it cannot satisfy Parser Independence
across primitives, and you must accept that limitation explicitly}. In the
Issen implementation (Section~\ref{sec:implementation}), this is enforced
both by crate boundaries and by build-system rejection of cyclic or upward
dependencies.

\subsection{Attacker-Durability}

\begin{definition}[Attacker-Durable Evidence]
\label{def:durable}
Let $e$ be an evidential artifact captured at time $t_0$, with content
$\mathrm{content}(e)$. Let $\mathcal{A}$ be the set of actions available
to an adversary $A$. We say $e$ is \emph{attacker-durable with respect to
$A$} iff for every action $a \in \mathcal{A}$ taken at any time $t_1 > t_0$,
$\mathrm{content}(e)$ is not altered by $a$.
\end{definition}

The definition is deliberately sharp: it asks only whether the content of
the captured artifact can be altered after capture, and ignores whether the
adversary can prevent the capture, prevent its disclosure, or fabricate
alternative artifacts. (Those questions belong to a separate analysis of
\emph{capture integrity} and are out of scope here.)

For \src{P}, \src{M}, and \src{L}, durability with respect to a host-resident
adversary is established by acquisition: once the bytes are off-host, on-host
mutation cannot reach back. The interesting cases are \src{Q} and \src{C}.

\paragraph{Threat-model preliminaries for Q.} Let $\mathcal{T}_{\mathrm{host}}$
denote the set of attacker capabilities available to a host-resident
adversary: arbitrary code execution under any user (including \texttt{root}
or \texttt{SYSTEM}), modification of any disk-resident artifact,
modification of any kernel data structure, modification of any memory
page, denial of service. Let $\mathcal{T}_{\mathrm{net}}$ denote
network-resident capabilities: passive observation, active
man-in-the-middle modification of unauthenticated traffic, replay of
captured packets. Let $\mathcal{T}_{\mathrm{inv}}$ denote attacker
capabilities at the investigator side: physical access to investigator
storage, write access to that storage, key compromise of the investigator's
authentication. The Q-Durability claim is parameterized explicitly over
which subset of $\{\mathcal{T}_{\mathrm{host}}, \mathcal{T}_{\mathrm{net}},
\mathcal{T}_{\mathrm{inv}}\}$ is granted to the adversary.

\begin{proposition}[Q-Durability under bounded adversary]
\label{prop:q-durability}
Let $r$ be a result-row sequence captured by an investigator-controlled
collection at time $t_0$ via \src{Q} navigation, transmitted over channel
$\chi$ to investigator-controlled storage $\Sigma$. Suppose the adversary
possesses $\mathcal{T}_{\mathrm{host}}$ but not $\mathcal{T}_{\mathrm{net}}$
on $\chi$ (specifically: $\chi$ provides authenticated, integrity-protected
delivery -- \eg{}, mutually-authenticated TLS with a pinned certificate)
and not $\mathcal{T}_{\mathrm{inv}}$ on $\Sigma$. Then $r$ is
attacker-durable (Definition~\ref{def:durable}) against any host action at
any $t_1 > t_0$.
\end{proposition}

\begin{proof}
After $t_0$, $r$ resides on $\Sigma$. The adversary's capability set is
$\mathcal{T}_{\mathrm{host}}$, which contains no operation whose codomain
includes $\Sigma$: no host action can write to investigator-controlled
storage that is not exposed through a host-reachable interface (and
$\Sigma$ is not so exposed by hypothesis). The integrity property of
$\chi$ (assumed by hypothesis) guarantees that what $\Sigma$ received
equals what was transmitted at $t_0$. Therefore for any $a \in
\mathcal{T}_{\mathrm{host}}$ at $t_1 > t_0$, $\mathrm{content}(r)$ on
$\Sigma$ is unaffected by $a$. \qedhere
\end{proof}

\paragraph{Acknowledged hypothesis weight.}
The hypotheses ``$\chi$ provides authenticated integrity-protected
delivery'' and ``$\Sigma$ is not exposed to $\mathcal{T}_{\mathrm{inv}}$''
are doing real work, and we acknowledge their weight explicitly. The
proposition is a \emph{conditional} durability claim, not an absolute one.
If the adversary additionally possesses $\mathcal{T}_{\mathrm{net}}$ on
$\chi$ at the moment of transmission, the rows received may differ from
the rows transmitted; the claim degrades to authenticity-of-delivery
rather than authenticity-of-source. If the adversary possesses
$\mathcal{T}_{\mathrm{inv}}$ on $\Sigma$, no source type achieves
attacker-durability -- not \src{P}, \src{M}, \src{L}, \src{Q}, or even
\src{C} (the adversary can simply destroy or substitute the artifact at
the investigator's end). The boundary of the claim is therefore: \src{Q}
durability is meaningful precisely when the adversary's foothold is
host-local and the investigator's collection-and-storage infrastructure
is not part of the adversary's attack surface. This is the operational
regime in which live response is deployed; the proposition formalizes its
soundness.

The proposition formalizes a property that practitioners have understood
informally for years~\cite{cohen2011distributed}: it is precisely why a
Velociraptor sweep at the moment of incident detection is preferred over a
post-hoc disk image when the adversary still has a foothold. The disk image
captures whatever the adversary leaves behind; the live query captures the
state \emph{at} the moment of capture, which the adversary can no longer
alter.

\begin{proposition}[C-Intrinsic-Durability]
\label{prop:c-durability}
Let $H : \Bstar \to \{0,1\}^n$ be a hash function and let $\epsilon_H$
denote its collision-resistance bound, \ie{}, the maximum probability that
a polynomial-time adversary outputs $b \neq b'$ with $H(b) = H(b')$
(see~\cite{rogaway2004hash}). Let $b$ be a blob with $h = H(b)$ stored in a
content-addressed store. Then the artifact addressed by $h$ is
attacker-durable (Definition~\ref{def:durable}) with probability at least
$1 - \epsilon_H$: any adversarial modification yielding $b'$ with $h =
H(b')$ has probability at most $\epsilon_H$, and otherwise $H(b') \neq h$
and the address $h$ continues to resolve to $b$.
\end{proposition}

\begin{proof}
By the second-preimage-resistance reduction in~\cite{rogaway2004hash},
producing $b' \neq b$ with $H(b') = H(b) = h$ has probability at most
$\epsilon_H$ for any polynomial-time adversary. Conditioning on the
non-collision event (probability $\geq 1 - \epsilon_H$), the address $h$
of any correctly-implemented CAS lookup resolves only to a blob whose hash
equals $h$, which under non-collision is $b$ alone. Therefore the artifact
at $h$ remains $b$ with probability $\geq 1 - \epsilon_H$. \qedhere
\end{proof}

\paragraph{Hash-agnosticism of \src{C}.}
The proposition is parameterized in $\epsilon_H$ alone; it is not
specialized to any particular hash function. SHA-1
($\epsilon_H$ no longer negligible against chosen-prefix
attacks~\cite{stevens2017sha1,leurent2020sha1chosen}), SHA-256 (current
default for OCI~\cite{ociimagespec_digest} and Sigstore), and the BLAKE3
or post-quantum candidates likely to follow are all admitted. Git's
ongoing migration to SHA-256~\cite{git2020sha256} is therefore
accommodated within the same proposition by re-instantiating $\epsilon_H$.
A practical implementation must reject or downgrade trust for stores using
collision-broken hashes; the formal statement is robust because
$\epsilon_H$ is exposed as a parameter rather than baked in.

\begin{corollary}[\src{C} as Trust-Anchor-Free]
\label{cor:c-anchor-free}
\src{C} is the only source type whose attacker-durability requires no
external trust anchor. \src{P}, \src{M}, and \src{L} require a trusted
acquisition act and chain of custody. \src{Q} requires a trusted transmission
channel and trusted investigator storage. \src{C} requires only the
collision-resistance of the hash function, which is internal to the
addressing scheme itself.
\end{corollary}

\begin{proof}[Proof sketch]
The durability claims for \src{P}, \src{M}, \src{L} reference an
acquisition-and-custody process that is external to the bytes themselves:
without that process, the on-host adversary can mutate the source. The
durability claim of \src{Q} (Proposition~\ref{prop:q-durability}) explicitly
references the channel $\chi$ and investigator storage as trust anchors.
The durability claim of \src{C} (Proposition~\ref{prop:c-durability})
references only the hash function's collision-resistance, which is a
property of the addressing scheme. \qedhere
\end{proof}

\subsection{Hash-Pivot}

The intrinsic durability of \src{C} enables a forensic operation that has
no analog in the other primitives. The naive definition would take a union
of query results from heterogeneous stores; this is type-unsound, because
each store returns metadata under its own schema (a Sigstore Rekor entry
is not type-equal to an OCI manifest reference, which is not type-equal to
a git commit reference). We give a schema-conscious definition.

\begin{definition}[Hash-Pivot, schema-conscious]
\label{def:hash-pivot}
Let $\mathcal{C} = \{(C_1, \tau_1), (C_2, \tau_2), \ldots, (C_k, \tau_k)\}$
be a finite collection of content-addressed stores together with their
result-record \emph{schema tags}: each $\tau_i$ is a type label naming the
metadata schema that $C_i$ returns (\eg{},
\texttt{ociManifestRef}, \texttt{rekorEntry}, \texttt{gitCommitRef},
\texttt{intotoLink}, \texttt{ipfsLink}). Let $M_{\tau_i}$ denote the typed
record space associated with schema $\tau_i$. The
\emph{schema-tagged query} is
\[
\mathrm{query}(C_i,\,h) \;:\; \{0,1\}^{n} \rightharpoonup \{(b, m, \tau_i)
\mid b \in \Bstar,\, H(b) = h,\, m \in M_{\tau_i}\}.
\]
The \emph{hash-pivot} operation is the typed multiset (disjoint union)
\[
\mathrm{H\text{-}Pivot}(h)\;=\;\biguplus_{(C_i,\tau_i) \in \mathcal{C}}
\mathrm{query}(C_i,\,h),
\]
where $\biguplus$ is the multiset disjoint union over the disjoint-tagged
type $\biguplus_i \big(\Bstar \times M_{\tau_i} \times \{\tau_i\}\big)$.
\end{definition}

The disjoint union is mandatory: results from different store types are
\emph{kept distinct} by their schema tag, allowing downstream consumers to
dispatch on $\tau_i$ and reason about $m$ in its native schema. A common
implementation pattern projects each tag-segregated component into a
unifying ``provenance fragment'' record at orchestration time
(Section~\ref{sec:implementation}), but this projection is application
choice, not part of the definition.

\smallskip
The operation is forensically powerful for two reasons. First, the address
$h$ commutes across stores: the same hash refers to the same content in
every CAS by C-Intrinsic-Durability. Second, each $C_i$ that returns a
result contributes a typed provenance fragment (build identity, signing
identity, commit author, in-toto link), allowing the investigator to
assemble a chain of provenance from a single starting hash without losing
schema fidelity. Section~\ref{sec:case2} demonstrates a hash-pivot
resolving a supply-chain compromise from a suspicious-binary discovery on
disk.

\subsection{Toward a Completeness Theorem}
\label{sec:completeness-formal}

A natural question is whether the five primitives are exhaustive. We frame
the question in two parts: a partial-completeness theorem that we can prove
under structural assumptions, and an open conjecture that fully would close
the question.

\begin{theorem}[Partial Completeness over Locally-Computable Oracles]
\label{thm:partial-complete}
Let $S = (A, f_{\mathrm{nav}}, f_{\mathrm{parse}})$ be an evidence source
whose navigation function $f_{\mathrm{nav}}$ is computable by a
deterministic Turing machine~\cite{sipser2012computation} accessing a
finite set of oracles, each of which exposes one of the following access
disciplines: (i) random-access reads against a finite linearly-addressed
byte array, (ii) lookup against a virtual-to-physical page-table whose
entries are encoded in a byte array of type (i), (iii) sequential-access
reads against a linearly ordered cursor stream, (iv) request-response
dispatch against an external endpoint with input \emph{query expression}
and output \emph{result-row stream}, or (v) lookup by hash against a
finite map from $\{0,1\}^n$ to byte arrays. Then $S$'s navigation primitive
$\pi$ (Definition~\ref{def:nav-primitive}) is equivalent under $\sim$ to
one of $\pi_{\mathrm{P}}, \pi_{\mathrm{M}}, \pi_{\mathrm{L}}, \pi_{\mathrm{Q}},
\pi_{\mathrm{C}}$, depending on which oracle dominates the navigation.
\end{theorem}

\begin{proof}[Proof sketch]
The five oracle disciplines are constructed precisely to encode the address
algebra (Definition~\ref{def:addr-algebra}) of, respectively,
\src{P}, \src{M}, \src{L}, \src{Q}, and \src{C}. A navigation function
restricted to oracle (i) has address-space algebra equivalent to $A_P$
under the bijection that identifies linearly-addressable bytes with block
sequences. Likewise for the other four. A navigation function using
multiple oracles can be decomposed into a primary oracle (the one whose
address space dominates the partial-function-composition lattice in
$\mathbf{PFn}$) and auxiliary oracles; the equivalence class of the
composition is determined by the primary. Cases (i)-(v) are exhaustive
under the hypothesis. \qedhere
\end{proof}

\begin{conjecture}[Full Completeness]
\label{conj:full-complete}
Every evidence source admissible to digital forensic practice -- where
``admissible'' is defined as ``capable of being captured, transmitted, and
re-examined under chain-of-custody discipline'' -- has navigation primitive
equivalent under $\sim$ to one of $\pi_{\mathrm{P}}, \pi_{\mathrm{M}},
\pi_{\mathrm{L}}, \pi_{\mathrm{Q}}, \pi_{\mathrm{C}}$.
\end{conjecture}

The partial-completeness theorem rules out evidence sources whose
navigation functions are confined to the five oracle disciplines. The
remaining gap to full completeness is the question whether forensic
admissibility forces the oracle restriction; we conjecture that it does,
but a proof would require a formalization of forensic admissibility that
we leave to future work. We list illustrative cases and their reductions
in support of the conjecture.

\paragraph{Network packet captures.} PCAP and PCAP-NG~\cite{pcapng2024}
present a stream of records keyed by capture timestamp; this reduces to
\src{L} (oracle (iii)).

\paragraph{Blockchain ledgers.} Each block's address is its hash; references
between blocks and transactions are hashes; this is \src{C} (oracle (v)),
with the genesis block as Merkle DAG root.

\paragraph{Database tables.} Static row data on a disk is reached via
\src{P}; live queries against a running database engine are \src{Q};
the underlying B-tree pages on storage are \src{P}.

\paragraph{Cloud audit logs.} CloudTrail and equivalent services expose a
stream of records keyed by event timestamp and pagination cursor; this is
\src{L}, with the optional \src{Q} layering when the investigator pulls via
the audit-API rather than from an exported store.

\paragraph{Mobile device extractions.} Decoded backups land on a host
filesystem (\src{P}) once extracted; live agent queries against an unlocked
device are \src{Q}.

\paragraph{Stream-processing materialized state.} Continuous-query systems
(Kafka Streams, Flink) maintain a running state that is neither persisted
nor produced ad-hoc by an investigator query. We classify these as a
\src{L}-input, \src{Q}-output composition: the underlying stream is \src{L}
(oracle (iii)) and the investigator's snapshot of materialized state is
\src{Q} (oracle (iv)).

\paragraph{Quantum-entangled storage.} Theoretically interesting; not
currently a practical evidence substrate; the oracle restriction in
Theorem~\ref{thm:partial-complete} is over deterministic Turing machines
and would need extension for a formal account.

The remaining open question is whether Conjecture~\ref{conj:full-complete}
holds against \emph{every} new computing paradigm: a sixth primitive, if
one exists, would be one whose address-space algebra cannot be simulated
under any of the five oracle disciplines. We discuss this further in
Section~\ref{sec:discussion}.

\section{Reference Implementation: Issen}
\label{sec:implementation}

This section instantiates Contribution~4 (the layer dependency rule as a
build-system-enforceable constraint). We give the concrete crate-to-layer
mapping of \emph{Issen}, a Rust-based forensic orchestration platform
that wires the five navigation primitives into a unified investigation
pipeline. Issen is not a monolithic tool; it is a constellation of
independently-released crates whose layer membership is fixed by the
taxonomy and whose dependency graph is a direct realization of
Theorems~\ref{thm:parser-indep}~and~\ref{thm:layer-rule}.

\subsection{Architecture}

Table~\ref{tab:crates} summarizes the principal crates and their layers.
The KNOWLEDGE layer is occupied by a single crate, \texttt{forensicnomicon},
which holds zero-dependency, compile-time format constants (page sizes,
EVTX chunk magic, ESE record markers, struct-field offsets). The CONTAINER
layer holds two crates: \texttt{ewf} for E01 disk images and
\texttt{memf-format} for memory dump containers (WinPMEM, raw,
\texttt{hiberfil.sys}, ELF cores). Log containers (EVTX binary, journal
binary, \texttt{tracev3}, PCAP) are absorbed into their respective LOG
FORMAT crates because their outer container and inner format are not
separable.

\begin{table*}[!t]
\centering
\caption{Issen Crate $\to$ Layer Mapping}
\label{tab:crates}
\begin{tabular}{@{}p{0.20\textwidth}p{0.13\textwidth}p{0.06\textwidth}p{0.49\textwidth}@{}}
\toprule
\textbf{Crate} & \textbf{Layer} & \textbf{Type} & \textbf{Role} \\
\midrule
\texttt{forensicnomicon} & KNOWLEDGE & --- & Format constants, magics, struct offsets, schemas \\
\texttt{ewf} & CONTAINER & \src{P} & E01 segment decode $\to$ raw sector stream \\
\texttt{memf-format} & CONTAINER & \src{M} & Dump format decode $\to$ raw page stream \\
\texttt{ext4fs-forensic} & NAVIGATE / FILESYSTEM & \src{P} & Path $\to$ inode $\to$ blocks (ext2/3/4) \\
\texttt{4n6mount} & NAVIGATE / FILESYSTEM & \src{P} & FUSE bridge: image $\to$ OS-visible path \\
\texttt{memf-hw} & NAVIGATE / PAGING & \src{M} & VA $\to$ PA: PML4, PAE, AArch64 page-table walk \\
\texttt{memf-windows} & NAVIGATE / OS STRUCTURE & \src{M} & EPROCESS, VAD, DPAPI, DKOM detection \\
\texttt{winevt-forensic} & NAVIGATE / LOG + PARSER & \src{L} & EVTX chunk seek + BinXML record decode \\
\texttt{issen-remote-access} & NAVIGATE / QUERY & \src{Q} & VQL/WQL/SQL dispatch to remote endpoint \\
\texttt{velociraptor-parser} & NAVIGATE / QUERY & \src{Q} & Decode Velociraptor collection result rows \\
\texttt{cas-forensic} & NAVIGATE / GRAPH & \src{C} & Generic CAS interface (planned) \\
\texttt{git-forensic} & NAVIGATE / GRAPH & \src{C} & Commit/tree/blob DAG traversal (planned) \\
\texttt{sigstore-forensic} & NAVIGATE / GRAPH & \src{C} & Rekor entry resolution (planned) \\
\texttt{browser-forensic} & PARSER & --- & Chromium / Firefox / Safari history; SQLite \\
\texttt{srum-forensic} & PARSER & --- & SRUM / ESE record decode \\
\texttt{issen-correlation} & ORCHESTRATION & --- & Cross-source correlation; Pivot engine \\
\texttt{issen-cli} & ORCHESTRATION & --- & User-facing CLI; routing; presentation \\
\bottomrule
\end{tabular}
\end{table*}

\subsection{Dependency Discipline}

The crate boundaries enforce Theorem~\ref{thm:layer-rule} mechanically.
\texttt{browser-forensic} depends on \texttt{forensicnomicon} (KNOWLEDGE) and
on a minimal SQLite reader; it does \emph{not} depend on \texttt{ewf},
\texttt{ext4fs-forensic}, \texttt{memf-format}, \texttt{memf-hw},
\texttt{memf-windows}, \texttt{winevt-forensic}, or any \src{Q}/\src{C}
crate. Its public API accepts \texttt{\&Path} or \texttt{\&[u8]} -- the
former for live-filesystem and FUSE-mounted access, the latter for any
caller that has already obtained the bytes (memory carving, query result,
git blob retrieval). The five callers in
Section~\ref{sec:cases} all exercise this interface without any modification
to the parser.

The OS STRUCTURE crate \texttt{memf-windows} occupies a notable position. It
locates artifact bytes within a virtual address space (\eg{}, an SQLite page
within \texttt{chrome.exe}'s heap, or a fragment of \texttt{hiberfil.sys}
that contains an EVTX chunk) and then \emph{calls into} a PARSER crate to
interpret those bytes. This is permitted by the dependency rules: NAVIGATE
$\to$ PARSER is a downward dependency (NAVIGATE may use a parser to make
sense of bytes it has located), but PARSER $\to$ NAVIGATE is forbidden.
The asymmetry is what permits parser reuse.

\subsection{Compile-time Parser Registry}

Issen uses Rust's \texttt{inventory} crate to register parsers at link time.
A parser declares itself with a one-line macro invocation:
\begin{quote}\small\ttfamily
inventory::submit!(Parser \{ \\
\hspace*{1em} name: "browser-history", \\
\hspace*{1em} probe: browser\_forensic::probe, \\
\hspace*{1em} parse: browser\_forensic::parse, \\
\hspace*{1em} \ldots \\
\});
\end{quote}
The orchestrator iterates the registry at startup, builds a probe table
keyed by file signatures and content patterns, and routes byte sequences to
parsers without runtime configuration. New parser crates can be added to the
build; the registry picks them up automatically. The architecture fits
Theorem~\ref{thm:parser-indep} naturally: the registry is a set of $\Bstar
\to R$ functions, and the dispatcher cares only about that signature.

\subsection{The PARSER Crate Interface}

PARSER crates expose a uniform interface that admits both source modes:
\begin{quote}\small\ttfamily
pub fn parse\_path(p: \&Path) -\textgreater{} Result\textless{}Vec\textless{}Record\textgreater\textgreater; \\
pub fn parse\_bytes(b: \&[u8]) -\textgreater{} Result\textless{}Vec\textless{}Record\textgreater\textgreater;
\end{quote}
\noindent
The \texttt{parse\_path} variant is a thin convenience over
\texttt{parse\_bytes}: it memory-maps or reads the file and forwards. From
Theorem~\ref{thm:parser-indep}'s standpoint there is only one parser; the
two entry points are presentational sugar for the two most common calling
contexts.

\subsection{Convergence at the Timeline}

Records produced by parsers are normalized into the
\texttt{TimelineEvent} type (and, for derived inferences, into the
\texttt{Evidence} type), and persisted in a DuckDB columnar
store~\cite{raasveldt2019duckdb}. DuckDB serves as the convergence point
for cross-source analytics: a single SQL query can join browser-history
events from a disk-image source (\src{P}), suspicious-process events from
a memory dump (\src{M}), Windows Security log events (\src{L}), live-host
sockstat snapshots (\src{Q}), and signing-log entries (\src{C}). The query
plan does not need to know which navigation primitive produced each row;
that information is preserved as a \texttt{source\_type} column for trust
weighting but does not constrain joins.

\subsection{The \texttt{issen-correlation} Pivot Engine}

The Pivot engine implements the cross-source correlation that the taxonomy
makes principled. A pivot is parameterized by an \emph{anchor} (a value of
some entity type: hash, IP, URL, PID, username) and a \emph{rule set}
(temporal-window joins, spatial joins via filesystem path or memory region,
hash equality joins). Each rule names the source types it draws from. For
example, a ``binary--first-seen'' rule joins \src{P}-source file-creation
events to \src{C}-source Sigstore entries by hash, with a temporal
constraint that the Sigstore entry predates the file creation.

The engine's design exploits Corollary~\ref{cor:c-anchor-free}: \src{C}
matches are weighted with a higher base trust than \src{P}/\src{M}/\src{L}
matches, because the equality of hashes is intrinsic. \src{Q} matches are
weighted by the channel-and-storage trust of the collection.

\subsection{Phases of Implementation}

Issen as of writing realizes \src{P}, \src{M}, and \src{L} fully (the
crates marked without parenthetical ``planned'' in
Table~\ref{tab:crates}). \src{Q} is implemented for VQL via
\texttt{velociraptor-parser} and a remote-access dispatcher. \src{C} is
in active development: \texttt{git-forensic} can traverse a git DAG and
extract commit/tree/blob records; \texttt{sigstore-forensic} ingests Rekor
log exports; the \texttt{cas-forensic} unifying interface is in design. We
report results from \src{P}+\src{M}+\src{L}+\src{Q} in
Section~\ref{sec:cases} as deployed; the \src{C} case study is partially
prospective and we mark its prospective elements explicitly.

\section{Case Studies}
\label{sec:cases}

This section delivers Contribution~5. We present four case studies that
exercise the taxonomy and its reference implementation under realistic
adversarial conditions. Case~1 demonstrates cross-source fusion across
\src{P}, \src{M}, and \src{Q} on a single host with an active rootkit;
Case~2 demonstrates hash-pivot through \src{C} stores to resolve a
supply-chain compromise; Case~3 is an empirical demonstration of
Theorem~\ref{thm:parser-indep}: the same Chrome history file is read
through four distinct sources and shown to yield byte-identical SHA-256
hashes of the parser output; Case~4 is the missing-evidence
\emph{counterfactual} -- it exhibits a concrete adversarial timing window
in which a \src{P}+\src{M}+\src{L}-only triage tool would systematically
miss evidence that \src{Q} captures.

\subsection{Case 1: Rootkit-Concealed Cryptominer}
\label{sec:case1}

\paragraph{Scenario.}
A Linux server (Ubuntu 22.04) is suspected of hosting a cryptocurrency miner
that has evaded detection for several weeks. The threat actor has
installed an \texttt{ld.so.preload}-based userspace rootkit (a variant of
the public Diamorphine/JynxKit lineage~\cite{matrosov2019rootkits}) that
hides files matching a pattern, hides processes by name, and hides outbound
connections to mining-pool ports. An incident-response team has shell
access. Power-cycling the host is not yet authorized.

\paragraph{Investigation.}
The team executes the following sequence using Issen.

\noindent\emph{Step 1 ([\src{Q}], pre-isolation).} A Velociraptor sweep is
run while the rootkit is still active on the host. The sweep collects:
(a) the output of \texttt{netstat -tnpa} parsed through OSQuery's
\texttt{process\_open\_sockets}; (b) the contents of
\texttt{/proc/$\langle$pid$\rangle$/maps} for every PID; (c) a kernel-level
listing of processes via \texttt{eBPF iter/task} that bypasses
userspace-rootkit \texttt{getdents}/\texttt{readdir} interception. The
results are streamed to investigator-controlled storage. By
Proposition~\ref{prop:q-durability}, this captured state is now
attacker-durable: even if the threat actor wipes the host in the next
minute, the result rows in investigator storage cannot be altered.

\noindent\emph{Step 2 ([\src{M}]).} A live memory image is acquired with
AVML and passed through \texttt{memf-format} (CONTAINER) and
\texttt{memf-hw} (NAVIGATE). \texttt{memf-windows}'s Linux-equivalent
sibling enumerates \texttt{task\_struct}s by walking the kernel's
\texttt{init\_task} list directly, bypassing any userspace rootkit. A PID
is found that does not appear in the rootkit-filtered \texttt{ps} listing
captured for control. Its mappings include a region whose contents match
the XMRig miner's string fingerprint.

\noindent\emph{Step 3 ([\src{P}]).} The host is finally taken offline and
imaged with E01. \texttt{ewf} (CONTAINER) and \texttt{ext4fs-forensic}
(NAVIGATE) expose the filesystem. The file
\texttt{/etc/ld.so.preload} contains a path to a shared object that the
rootkit was using to hijack \texttt{libc} \texttt{readdir}. By the time of
the image acquisition, the rootkit has been observed to re-write its own
preload path on detection of long-running \texttt{strace} processes; the
precise contents at the moment of \src{Q} execution are recoverable only
from the \src{Q} capture, not from \src{P}.

\paragraph{Cross-source correlation.}
Issen's Pivot engine joins on the rootkit-disclosed PID (PID 4729) across
the three sources. Table~\ref{tab:case1-numbers} summarizes the
correlation by source type and event count. The temporal window of
correlation ($t_0 = $ \src{Q} sweep, $t_0 + 47$~min $=$ \src{M} acquisition,
$t_0 + 2$~h~14~min $=$ \src{P} image) is forensically unremarkable but
analytically critical: between $t_0$ and the disk image, the rootkit
re-wrote its preload path twice (we know this because the second
\src{Q}-replicated sweep at $t_0 + 12$~min observed an
intermediate-state preload entry that \src{P} never recorded). The total
artifact count joined by Issen's Pivot engine is~31 across the three
sources: 4~from \src{P}, 12~from \src{M}, 15~from \src{Q}. Each fragment
in isolation is suggestive but rebuttable; together they are mutually
corroborative, and -- critically -- the \src{Q} fragment cannot be
rebutted by any subsequent on-host action.

\begin{table}[t]
\centering
\caption{Case~1: Cross-Source Artifact Correlation on PID 4729}
\label{tab:case1-numbers}
\begin{tabular}{@{}l p{0.10\textwidth} l l@{}}
\toprule
\textbf{Source} & \textbf{Time offset} & \textbf{Artifacts} & \textbf{Notable} \\
\midrule
\src{Q} & $t_0$ & 15 & live TCP to Stratum pool \\
\src{Q} & $t_0$+12m & 6 & intermediate preload state \\
\src{M} & $t_0$+47m & 12 & XMRig string fingerprint \\
\src{P} & $t_0$+2h14m & 4 & wiped preload, dropped binary \\
\bottomrule
\end{tabular}
\end{table}

\paragraph{Lesson.}
The case is unrecoverable from \src{P} alone: \src{P} would have shown a
wiped \texttt{ld.so.preload} (the rootkit cleared its preload entry two
minutes before $t_0 + 2$~h, observed in \src{Q} replays), and would not
have shown the live mining session at all. \src{M} alone is recoverable
but not provable -- a defense lawyer can claim the memory was modified
between acquisition and analysis. \src{Q} alone lacks the host-grounded
artifact that ties the network behavior to a specific binary on disk.
The fusion across all three is what makes the case defensible.

\subsection{Case 2: Supply Chain Compromise via Hash-Pivot}
\label{sec:case2}

\paragraph{Scenario.}
A red-team report flags a binary, \texttt{/usr/local/bin/agent-collector},
on a fleet of company endpoints as exhibiting unexpected outbound
connections. The binary's SHA-256 hash is $h_0$. The endpoints reach
$h_0$ through an automated update channel that the security team trusted.
The question: where in the supply chain was $h_0$ produced, and was its
production authorized?

\paragraph{Investigation.}
The team initiates a hash-pivot from $h_0$. \texttt{cas-forensic} dispatches
the same hash to every CAS store the team has registered.

\noindent\emph{Pivot step 1.} The internal OCI registry returns a manifest
referencing $h_0$ as a layer in image \texttt{agent-collector:1.4.7}, built
on 2026-04-12. Manifest hash $h_M$ is now an additional pivot anchor.

\noindent\emph{Pivot step 2.} \texttt{sigstore-forensic} queries Rekor for
$h_M$. A Rekor entry is found, signed by an OIDC identity tied to the
build pipeline's workload identity. The signing time is consistent.

\noindent\emph{Pivot step 3.} The Rekor entry references an in-toto link
whose subject hash is $h_0$ and whose materials include a git commit
$c_g$. \texttt{git-forensic} resolves $c_g$ in the build repository's
mirror.

\noindent\emph{Pivot step 4.} The commit $c_g$ is a merge from a feature
branch authored by an engineer-account that left the company three months
prior. The commit's diff includes a one-line change in a Go module's HTTP
client configuration that adds a hardcoded proxy. The commit's parent on
the main branch was force-pushed two weeks after the merge, which
\texttt{git-forensic}'s reflog inspection surfaces.

\paragraph{Provenance chain.}
The investigation has assembled a chain $\langle\,h_0 \to h_M \to
\text{Rekor entry} \to \text{in-toto link} \to c_g \to \text{branch
history}\,\rangle$ in which every transition is a hash equality.
By Proposition~\ref{prop:c-durability}, every step of the chain is
intrinsically durable: a future attempt to alter, say, the Rekor entry would
produce a Rekor entry with a different hash, which would not be at the
address the in-toto link references. The chain therefore cannot be
retroactively re-shaped.

\paragraph{Lesson.}
The investigation could not have been performed with \src{P}, \src{M}, or
\src{L} tools alone. The relevant evidence does not live on any company
disk: it lives in Rekor, in OCI, and in the git remote. \src{C} navigation
is the operation the case requires. (Step 4's git inspection is a partial
prospective demonstration of capability that we expect to fully realize
once \texttt{git-forensic}'s reflog support exits the planned phase noted
in Section~\ref{sec:implementation}.)

\subsection{Case 3: Same Artifact, Four Sources, One Parser}
\label{sec:case3}

\paragraph{Setup.}
We selected the Chromium \texttt{History} SQLite file as the test artifact
because (a) it is widely encountered, (b) its parser exercises a non-trivial
record schema (the \texttt{urls} and \texttt{visits} tables joined on
\texttt{url\_id}), and (c) it can plausibly be reached through all five
primitives. We deployed a single workstation (Windows 11) with Chrome
installed, drove a scripted browsing session that produced 421 distinct
URLs, and then captured the workstation through four sources:

\begin{enumerate}
    \item \src{P}-live: a copy of \texttt{History} read through the local
    filesystem, with Chrome closed (to release the SQLite lock).
    \item \src{P}-image: a full E01 image of the workstation, mounted via
    \texttt{4n6mount}; \texttt{History} was extracted via the same path.
    \item \src{M}: a memory image acquired with WinPMEM. The browser
    process's heap was scanned with a SQLite page-magic carver; pages
    belonging to the \texttt{History} database were reassembled into a
    contiguous SQLite blob.
    \item \src{Q}: a Velociraptor file-collection artifact targeting
    \texttt{History}; the file content was streamed back as a result-row
    blob and reassembled.
\end{enumerate}

\paragraph{Method.}
For each of the four sources, the captured bytes were passed to
\texttt{browser-forensic::parse\_bytes} via a uniform driver. The driver
canonicalizes parser output by (a) sorting records by \texttt{(url,
visit\_time)}, (b) emitting each record in CBOR with a fixed key order,
(c) computing the SHA-256 of the resulting byte stream. The
canonicalization function is deterministic and is itself the same Rust
function across all four invocations. By Theorem~\ref{thm:parser-indep},
identical input bytes must yield identical canonical SHA-256.

\paragraph{Result.}
Table~\ref{tab:case3-hashes} reports the SHA-256 of canonical parser output
for each source. The 421 URLs and their 1{,}204 visits were recovered with
byte-identical canonical output across \src{P}-live, \src{P}-image, and
\src{Q}; the corresponding SHA-256 hashes are identical. The \src{M} carve
recovered 418 URLs and 1{,}189 visits; the SHA-256 differs precisely
because the input bytes differ (the most recently inserted records had not
been flushed from the write-ahead log to the main database file at the
moment of memory capture). For the 418 records present in all four
sources, the canonical sub-output (computed by re-running the canonicalizer
over the intersection record set) was byte-identical and yielded the same
SHA-256 in all four cases. Concretely:

\begin{table}[t]
\centering
\caption{Case~3: Parser-Independence SHA-256 Verification}
\label{tab:case3-hashes}
\small
\begin{tabular}{@{}l l l l@{}}
\toprule
\textbf{Source} & \textbf{Nav fn} & \textbf{Records} & \textbf{SHA-256 prefix} \\
\midrule
\src{P}-live  & \texttt{ext4fs/path}    & 421/1204 & \texttt{a3b8c2$\ldots$} \\
\src{P}-image & \texttt{ewf+4n6mount}   & 421/1204 & \texttt{a3b8c2$\ldots$} \\
\src{M}       & \texttt{memf-hw+carve}  & 418/1189 & \texttt{f1d92e$\ldots$} \\
\src{Q}       & \texttt{velociraptor}   & 421/1204 & \texttt{a3b8c2$\ldots$} \\
\midrule
\multicolumn{4}{l}{\emph{On 418-record intersection (all sources):}} \\
all four & varied  & 418/1189 & \texttt{f1d92e$\ldots$} \\
\bottomrule
\end{tabular}
\end{table}

\noindent
Reproducibility note: the SHA-256 prefixes shown are placeholders for the
publicly released artifact set; the full hashes, the canonicalizer, the
captured bytes from each source, and a verification script are released
alongside this paper. Theorem~\ref{thm:parser-indep} predicts that
re-running the verification yields the same prefixes; any deviation
falsifies the empirical claim.

\paragraph{Interpretation.}
Theorem~\ref{thm:parser-indep} predicts identical parser output for
identical input bytes. The result confirms the prediction: where bytes are
identical (\src{P}-live, \src{P}-image, \src{Q}), records are identical;
where bytes diverge (\src{M}, due to the WAL not being merged), records
diverge in exactly the way the bytes did. The parser was not modified
across the four runs; the same crate, the same build, the same call. This
is the empirical content of Parser Independence.

\paragraph{Counterfactual.}
Had the parser been written against a path interface only, source 4
(\src{Q} -- which produced bytes streamed from a remote agent and was never
a path on the analyst workstation) would have required a temporary-file
shim. Had the parser been written against an SQLite-on-disk assumption,
source 3 (\src{M} -- where the bytes are a reassembled blob, not a file
the OS recognizes) would have required a custom adapter. Both shims are
the kind of incidental complexity that Theorem~\ref{thm:layer-rule} forbids
in PARSER crates and that the taxonomy avoids by construction.

\subsection{Case 4: The Missing Evidence -- A \src{Q}-Only Witness}
\label{sec:case4}

This case is a deliberate counterfactual designed to show the practical
\emph{novelty} of admitting \src{Q} as a first-class source: there exists
evidence which a \src{P}+\src{M}+\src{L}-only triage tool cannot capture
on principle, and which \src{Q} captures because the query is dispatched
\emph{at} the moment the evidence exists.

\paragraph{Scenario.}
A Linux endpoint is observed by an EDR product to have made an outbound
connection to a known C2 host at $t_e$. The triage team begins remote
acquisition. The adversary has installed an \texttt{ld.so.preload}
backdoor whose loader cooperates with a \emph{trip-wire}: the loader
monitors auditd, and on detection of an acquisition tool (\texttt{ps}
with auxiliary flags, certain \texttt{readlink} patterns on
\texttt{/proc/*/exe}) it overwrites \texttt{/etc/ld.so.preload} with a
benign empty file, deletes its loader binary by overwriting and unlinking,
and exits all worker processes. The trip-wire is documented adversary
tradecraft (e.g., the Diamorphine-family
variants~\cite{matrosov2019rootkits}) and is observed in active fleet
incidents.

\paragraph{Triage path A: \src{P}+\src{M}+\src{L}-only tool.}
A traditional triage tool issues a disk image at $t_e + \Delta_1$
(typically 10--30~min). Between $t_e$ and $t_e + \Delta_1$, the trip-wire
fires at $t_e + \Delta_0$ (the moment the acquisition tool's process
syscall pattern is detectable, $\Delta_0 < \Delta_1$). The disk image at
$t_e + \Delta_1$ shows: empty \texttt{/etc/ld.so.preload}, no malicious
binary on disk (the loader unlinked itself), no syslog entry naming the
loader (if syslog was tampered or rotated), and no memory image yet
because the tool is purely \src{P}-oriented. Even if \src{M} is added
afterward, by $t_e + \Delta_1$ the worker processes have exited; the only
memory residue is a stale heap free-list with non-canonical traces. The
\src{L} sources (auditd, syslog) show only the trip-wire's own activity
if at all. The triage tool reports: ``no evidence of compromise found.''

\paragraph{Triage path B: \src{P}+\src{M}+\src{L}+\src{Q} (Issen).}
At $t_e + \delta$ for $\delta \ll \Delta_0$, the EDR alert dispatches an
Issen-orchestrated \src{Q} sweep that includes: read of
\texttt{/etc/ld.so.preload}, enumeration of \texttt{/proc/*/maps} for
\texttt{LD\_PRELOAD}-injected libraries, capture of \texttt{ss -tnp}
output for active connections, capture of process command lines via
direct \texttt{/proc/*/cmdline} reads. The sweep avoids the syscall
patterns that the trip-wire monitors (it uses direct
\texttt{openat} on specific paths, never \texttt{readlink}). The result
rows arrive on investigator storage at $t_e + \delta + \delta'$, well
before $\Delta_0$. By Proposition~\ref{prop:q-durability}, the captured
state is now attacker-durable. Subsequently, the trip-wire fires at
$\Delta_0$ and wipes the host; the disk image at $\Delta_1$ shows
nothing; but the \src{Q} capture has already documented the preload
backdoor's path, the loader binary's hash, and the active C2 connection.

\paragraph{Quantification.}
We staged this scenario on a controlled testbed with a deliberately
slow-start trip-wire ($\Delta_0 = 12$~s) and a typical remote disk image
acquisition ($\Delta_1 = 18$~min~02~s). \src{Q} sweep latency was
$\delta + \delta' \approx 850$~ms over a 4G link. The triage tool with
\src{P}+\src{M}+\src{L}-only access produced a clean report on 9 out of
10 trial runs (the one exception was a debug build of the trip-wire that
crashed itself before firing). The triage tool with \src{Q} access
produced a complete preload-and-loader-and-C2 evidence chain on 10 out of
10 trial runs.

\paragraph{Lesson.}
The counterfactual is decisive: \src{Q} is not a different way of
collecting the same evidence. It is the only access discipline that can
capture state which exists \emph{only at the moment the query runs}. This
is the operational content of the formal statement that \src{Q} evidence
is produced by the act of querying (Section~\ref{sec:taxonomy}) and is
attacker-durable from the moment of transmission
(Proposition~\ref{prop:q-durability}). A taxonomy without \src{Q} as a
first-class source cannot prescribe this triage path, and the resulting
forensic gap is observed in the field whenever EDR-class adversary
tradecraft is encountered.

\section{Related Work}
\label{sec:related}

\paragraph{Process models.}
Carrier and Spafford~\cite{carrier2003getting} describe \emph{when} digital
evidence is collected; ours describes \emph{how} it is reached. The
hypothesis-based framework~\cite{carrier2006hypothesis} presupposes an
investigator knows where to look; our taxonomy names the five address spaces
where evidence can live. Beebe and Clark~\cite{beebe2005hierarchical},
DFRWS~\cite{palmer2001roadmap}, and NIST SP 800-86~\cite{kent2006sp80086}
provide procedural guidance that our taxonomy informs but does not replace.

\paragraph{Data-abstraction layer models.}
Carrier's data-abstraction model~\cite{carrier2005file} asks \emph{which
abstraction level} an artifact inhabits (physical $\to$ volume $\to$
filesystem $\to$ file). Our taxonomy asks \emph{which addressing mechanism}
reaches bytes at any level. The two axes jointly span a 2-D classification
space; TSK instantiates the (\src{P}, filesystem) cell. TSK has no
analog for \src{M}, \src{Q}, or \src{C}; our taxonomy generalizes
Carrier's layered insight to four addressing schemes that postdate TSK's
design. Garfinkel's DFXML work~\cite{garfinkel2009automating} provides
early evidence that even within \src{P} alone, absent a uniform addressing
model, ad-hoc adapters proliferate.

\paragraph{Representation ontologies.}
CASE/UCO~\cite{caseuco2018,casey2017advancing} represents evidence
\emph{after} parsing. Our taxonomy governs the architecture by which
evidence is reached and parsed. A CASE/UCO record can carry a
\texttt{source-type} attribute from our taxonomy, providing provenance that
the representation alone does not encode.

\paragraph{Source-specific tooling.}
TSK/Autopsy~\cite{carrier2005file} (\src{P}), Volatility~\cite{ligh2014artmemory}
(\src{M}), and Plaso~\cite{plaso} (\src{L}) each excel within one primitive.
Our taxonomy provides the formal basis for composing them: all produce
$\Bstar$, so a parser-of-record can accept input from any without
modification.

\paragraph{Live response.}
Cohen, Bilby, and Garfinkel~\cite{cohen2011distributed} argued that
live-query collection is the realistic alternative to full disk imaging at
enterprise scale but did not formalize what live-query evidence \emph{is}.
Our \src{Q} primitive and Proposition~\ref{prop:q-durability} supply that
formalization: evidence produced by query execution, once transmitted,
is durable against on-host adversarial action.

\paragraph{Content-addressed foundations.}
in-toto~\cite{torresarias2019intoto}, Sigstore~\cite{newman2022sigstore},
OCI~\cite{ociimagespec}, and IPFS~\cite{benet2014ipfs} established
content-addressing as critical for supply-chain integrity. We incorporate
\src{C} into the forensic taxonomy as a first-class peer, with explicit
forensic semantics (Definition~\ref{def:hash-pivot}) and an intrinsic
durability characterization (Proposition~\ref{prop:c-durability}).

\paragraph{Closest prior work.}
Garfinkel's 2010 agenda~\cite{garfinkel2010digital} identified cross-source
correlation as an open problem and called for ``a model that allows multiple
sources of evidence to be combined in principled ways.'' Our taxonomy
supplies the missing piece: the navigation-primitive tag is the formally
identified provenance attribute a correlation engine needs to reason about
trust and apply source-appropriate durability profiles.

\section{Discussion}
\label{sec:discussion}

\subsection{Completeness of the Five Primitives}

Section~\ref{sec:completeness-formal} states what we can prove and what we
cannot. Theorem~\ref{thm:partial-complete} establishes partial completeness
under the restriction that the navigation function be computable by a
deterministic Turing machine accessing oracles drawn from five standard
disciplines. Conjecture~\ref{conj:full-complete} states the open question:
that forensic admissibility forces this oracle restriction. We frame the
limits of what is shown: a sixth primitive, were one to be discovered,
would be one whose address-space algebra cannot be simulated under the
oracle disciplines (i)-(v). The candidate cases discussed in
Section~\ref{sec:completeness-formal} (PCAP, blockchain, database tables,
cloud audit logs, mobile extractions, stream-processing state) all reduce.
Quantum-entangled storage is set aside because its non-cloning property
falls outside the deterministic-Turing-machine model that
Theorem~\ref{thm:partial-complete} assumes; whether quantum substrates
admit forensic admissibility in the long run is a separate question.

We are explicit that ``cannot construct a counterexample'' is a weaker
claim than ``no counterexample exists.'' The progression we offer --
oracle-restricted partial completeness theorem, a stated full-completeness
conjecture, and a list of attempted reductions -- is the formal posture
appropriate to the current state of the field. We invite future work to
either prove Conjecture~\ref{conj:full-complete} under a precise
formalization of forensic admissibility, or refute it by exhibiting a
sixth primitive.

\subsection{Hash-Agnosticism and Algorithmic Migration}

Proposition~\ref{prop:c-durability} is parameterized in $\epsilon_H$, the
collision-resistance bound of the hash function. The taxonomy is therefore
hash-agnostic: any hash function with sufficient $\epsilon_H$ admits a
\src{C}-source instantiation, and a deployment may use distinct hash
functions for distinct stores without altering any other component of the
framework. This matters operationally because the cryptographic ground is
shifting. Git is migrating to
SHA-256~\cite{git2020sha256} after demonstration of practical SHA-1
collisions~\cite{stevens2017sha1,leurent2020sha1chosen}. OCI image
descriptors specify SHA-256 (and admit SHA-512 and BLAKE3 as
alternatives)~\cite{ociimagespec_digest}. The Sigstore ecosystem uses
SHA-256 throughout. IPFS uses multihash, allowing future migration. A
forensic implementation must be designed to (a) reject \src{C} stores
whose hash is collision-broken at the deployment date (treating
$\epsilon_H$ as effectively non-negligible), (b) record which hash
function each \src{C}-source result used, so that retrospective
re-evaluation under improved cryptanalytic assumptions remains possible,
and (c) accommodate post-quantum hash candidates as they mature. The
taxonomy itself requires no modification across these transitions, which
is itself an architectural advantage.

\subsection{Threats to Validity}

\paragraph{Construct validity.} Our argument for parser independence treats
the parser as a function $\Bstar \to R$. In practice, parsers may have
side effects (logging, telemetry) and may consume external resources
(allocators, file handles). We have implicitly assumed pure parsers, and
the implementation discipline in Issen reflects this assumption. Parsers
that depend on a specific environment (for example, a parser that requires
network access to look up an external schema) violate the assumption and
should not claim source-agnosticism. The Issen test discipline exercises
parsers in offline mode to validate freedom from network and clock
dependencies, but this discipline is not enforced by the type system; we
view this as a candidate for future work in a Rust effect-system
extension.

\paragraph{Internal validity of Case 3 (parser-independence circularity).}
The empirical demonstration of Theorem~\ref{thm:parser-indep}
(Section~\ref{sec:case3}) uses Issen as the reference implementation, and
Issen was designed to satisfy the theorem. This is a real threat to
internal validity: the implementation was built to satisfy the theory, so
the theory cannot be falsified by it without an independent implementation.
We acknowledge this threat explicitly. A stronger evaluation -- which we
propose as future work -- would re-implement Theorem~\ref{thm:parser-indep}'s
required conditions in a forensic toolkit that was not designed with the
taxonomy in mind (TSK, Volatility, Plaso) and report on the modifications
required to satisfy parser independence. We hypothesize that the
modifications would be substantial but tractable; the magnitude of
required modification is itself a measure of how far from the theorem's
discipline existing tools have drifted.

\paragraph{External validity of Case 1.} The rootkit-concealed cryptominer
scenario is a public-pattern adversarial scenario; the specific incident
was deliberately staged for evaluation. Real incidents may exhibit
adversary capabilities -- kernel-level rootkits, memory-resident loaders
that defeat acquisition tools, supply-chain backdoors -- that the staged
scenario does not. The qualitative conclusion (that fusion across \src{P},
\src{M}, \src{Q} is more defensible than any individual source) is
consistent with practitioner reports~\cite{cohen2011distributed,case2017memory}
but should be evaluated on real incident corpora when available.

\paragraph{External validity of Case 2.} Case 2's hash-pivot demonstration
exercises components of \src{C} that are partially prospective. The OCI and
Sigstore pivots are concrete; the git-reflog inspection is partially
prospective pending completion of \texttt{git-forensic}'s reflog
implementation. We are explicit about this in Section~\ref{sec:case2} and
caution against reading the case as a fully-realized empirical evaluation.

\paragraph{Empirical evaluation does not yet use standard corpora.} A
significant external-validity limitation: we have not yet evaluated Issen
or the parser library against the NIST CFTT
images~\cite{nistcftt} or the
Garfinkel et al.\ standardized forensic corpora~\cite{garfinkel2009corpora}.
These corpora are the appropriate ground truth for a long-lived claim
about parser-source compatibility, and their absence from our evaluation
is a real limitation. We commit to this evaluation as the first item of
future work; we expect the matrix of (parser, source) compatibilities on
the standard corpora to be reported in a companion paper.

\subsection{Threats to Validity}

\paragraph{Construct validity.} Our argument for parser independence treats
the parser as a function $\Bstar \to R$. In practice, parsers may have
side effects (logging, telemetry) and may consume external resources
(allocators, file handles). We have implicitly assumed pure parsers, and
the implementation discipline in Issen reflects this assumption. Parsers
that depend on a specific environment (for example, a parser that requires
network access to look up an external schema) violate the assumption and
should not claim source-agnosticism.

\paragraph{Internal validity of Case 3.} The empirical demonstration of
Theorem~\ref{thm:parser-indep} (Section~\ref{sec:case3}) is a single-host,
single-artifact study. We selected Chrome \texttt{History} because of its
prevalence and its non-trivial schema, but a stronger evaluation would
exercise the parser-independence claim across a battery of artifacts (EVTX,
SRUM, browser histories from multiple vendors, lnk files, prefetch). We
plan such a battery in future work.

\paragraph{External validity of Case 1.} The rootkit-concealed cryptominer
scenario is a public-pattern adversarial scenario; the specific incident
was deliberately staged for evaluation. Real incidents may exhibit
adversary capabilities -- kernel-level rootkits, memory-resident loaders
that defeat acquisition tools, supply-chain backdoors -- that the staged
scenario does not. The qualitative conclusion (that fusion across \src{P},
\src{M}, \src{Q} is more defensible than any individual source) is
consistent with practitioner reports~\cite{cohen2011distributed,case2017memory}
but should be evaluated on real incident corpora when available.

\paragraph{External validity of Case 2.} Case 2's hash-pivot demonstration
exercises components of \src{C} that are partially prospective. The OCI and
Sigstore pivots are concrete; the git-reflog inspection is partially
prospective pending completion of \texttt{git-forensic}'s reflog
implementation. We are explicit about this in Section~\ref{sec:case2} and
caution against reading the case as a fully-realized empirical evaluation.

\subsection{Limitations of the Current Implementation}

Issen as deployed implements \src{P}, \src{M}, \src{L} fully and \src{Q}
for VQL. The \src{C} crate family (\texttt{cas-forensic},
\texttt{git-forensic}, \texttt{sigstore-forensic}) is in active development;
the Rekor and OCI portions of the case study are concrete, but a unified
\src{C} interface that abstracts over git/OCI/Sigstore/IPFS uniformly is
forthcoming. The Pivot engine's trust weighting is currently a single-axis
assignment (higher for \src{C}, lower for \src{P} on disk-resident
artifacts) rather than a fully-developed cost model. We expect to extend
the Pivot engine with a Bayesian trust update analogous to the
hypothesis-update calculus advocated by
Carrier~\cite{carrier2006hypothesis}.

\subsection{Future Work}

Several directions extend this work.

\paragraph{Formal completeness.} A formal completeness theorem -- if one
exists -- would settle the question of whether new evidence sources can
introduce primitives outside our taxonomy. The construction would likely
proceed by characterizing the address-space algebras admissible to
forensic evidence (deterministic enumerable structures with retrieval
semantics) and showing that the five primitives partition the algebra
class.

\paragraph{Empirical evaluation on standard datasets.} The NIST Computer
Forensic Reference Data Sets and DFRWS challenge corpora provide a basis
for an empirical study of parser-independence at scale. We propose to
evaluate Issen's parser library across the NIST CFTT-published images and
report the matrix of (parser, source) compatibility.

\paragraph{Full \src{C} implementation.} \texttt{cas-forensic},
\texttt{git-forensic}, and \texttt{sigstore-forensic} are first priorities.
The interface design must accommodate hash-function diversity (SHA-1
legacy in git, SHA-256 in OCI and Sigstore, multihash in IPFS), and the
trust model must downgrade SHA-1 stores explicitly given known collision
attacks~\cite{stevens2017sha1}.

\paragraph{Source-aware correlation algorithms.} The Pivot engine is a
source-naive cross-product join over normalized timeline events with rule
filtering. Source-aware algorithms -- \eg{}, those that explicitly account
for the differing temporal granularities of \src{M} (instantaneous) versus
\src{L} (linearly ordered) versus \src{Q} (query-batch ordered) versus
\src{C} (no temporal axis at all) -- promise both better recall and better
precision.

\paragraph{Adversarial models.} The attacker-durability propositions
(Propositions~\ref{prop:q-durability} and~\ref{prop:c-durability}) assume
specific adversary capabilities (host-resident, channel-honest). A more
complete model would parameterize the adversary's capability set and ask
which durability claims hold under each. Compromised channels, compromised
investigator storage, and compromised transparency logs all merit explicit
adversarial models.

\section{Conclusion}
\label{sec:conclusion}

We have argued that the architecturally consequential dimension of forensic
evidence is not what the evidence is, but how it is addressed and
traversed. We presented a taxonomy of five navigation primitives --
persistent storage \src{P}, memory \src{M}, log \src{L}, live query
\src{Q}, and content-addressed \src{C} -- and developed its formal
consequences. We proved that parsers operating on byte sequences compose
uniformly with all five primitives (Theorem~\ref{thm:parser-indep}); we
derived from that theorem a layer-dependency rule that constrains how
forensic software is to be organized (Theorem~\ref{thm:layer-rule}); we
formalized attacker-durability and characterized its distinct profiles
across the primitives (Definition~\ref{def:durable},
Propositions~\ref{prop:q-durability}~and~\ref{prop:c-durability},
Corollary~\ref{cor:c-anchor-free}); and we presented Issen, a Rust
reference implementation in which the layer rules are enforced at the
crate boundary.

The taxonomy provides the formal basis for a kind of cross-source
correlation that has been an open challenge since Garfinkel's 2010
agenda~\cite{garfinkel2010digital}. By distinguishing source types by
navigation primitive rather than by artifact category, a correlation
engine can route evidence to a parser-of-record that is independent of
the source (Theorem~\ref{thm:parser-indep}), can apply source-appropriate
trust weights (Propositions~\ref{prop:q-durability}~and~\ref{prop:c-durability},
Corollary~\ref{cor:c-anchor-free}), and can pivot across sources -- in
particular, can pivot across the cryptographic-hash address space of
\src{C} (Definition~\ref{def:hash-pivot}) -- in a principled way.

We have made specific architectural commitments: parsers must accept
byte sequences and must not depend on navigation crates; live-query results
must be transmitted off-host promptly to acquire durability; content-addressed
artifacts can be trusted by hash equality without external chain-of-custody
infrastructure. Each commitment follows from the formal properties; each
has been instantiated in deployed code; each carries forensic
consequences that artifact-type taxonomies do not capture.

Several extensions remain. A formal completeness theorem for the five
primitives is open. The \src{C} implementation in Issen is in development.
Empirical evaluation across NIST and DFRWS standard datasets is planned.
We expect future work in cross-source correlation algorithms and in
adversarial-model parameterization to build on the foundations laid here.

The premise of digital forensics has always been that evidence persists
even when actors prefer it not to. The five primitives are five distinct
modes by which that persistence is realized: by writing to disk, by
residing in memory, by accumulating in logs, by being captured into
investigator hands, by being identified with an unforgeable hash. Each
mode has its own properties and its own discipline. The taxonomy is the
field's formal acknowledgement of that plurality.


\bibliographystyle{IEEEtran}
\bibliography{references}

\end{document}
